% section_9_5_before.tex
% The original §9.5 — written by Claude in the first draft pass.
% This is the version with the AI tells in place. The "after" version
% (also in this folder) shows what the proofread-and-approve loop
% produced.

\subsection{Interpretation and caveats}
\label{subsec:interpretation}

The joint reading is clean. New trade theory's prediction --- that
composition is downstream of MA and carries no independent growth
signal --- holds for distinctiveness and entropy, as Marshall
externalities would predict against if anything. The
Romer-Jacobs/Jacobs-externality prediction --- that any breadth of
knowledge drives growth --- is supported at marginal significance by
the total-diversity coefficient. The Mokyr useful-knowledge prediction
is supported at conventional significance through two distinct margins:
the extensive margin (share-of-useful, driven by vernacular composition
shifts) and the intensive margin (within-useful diversity, driven by
Latin scholarly depth). The cross-language pattern is not noise; it is
the predicted pattern under a Mokyr channel that operates through
different print media for different audiences.

Caveats are in order. (i) The diversity and share-of-useful coefficients
are small in absolute magnitude --- a one-standard-deviation increase
in share-of-useful is associated with about a 1.5 log-point increase
in 50-year city growth. (ii) The growth channel documented here may
operate at a longer horizon than our 1450--1650 window suggests, and
our reduced-form panel cannot distinguish equilibrium-conditioned
composition coefficients from causal ones. (iii) The Thirty Years' War
(1618--1648) shapes the 1600--1650 growth window for German-territory
cities; Appendix~\ref{app:war} reports the sensitivity --- both §8
composition signals and §9 diversity-on-growth strengthen once
war-affected observations are removed. (iv) The cross-language pattern
is real and we do not claim a single channel explains every
coefficient: Italian print shows a significantly negative growth
coefficient on diversity ($-0.251^{**}$,
Table~\ref{tab:section9_main}) that runs against the pooled positive,
plausibly reflecting the outsized weight of Venezia in the
Italian-print panel (35\% of all Italian titles; see
Section~\ref{sec:geography}) over a period in which Venezia's own
economic position weakened.
